\documentclass[conference]{IEEEtran}
\IEEEoverridecommandlockouts

\usepackage{cite}
\usepackage{amsmath}
\usepackage{graphicx}
\usepackage{booktabs}
\usepackage{url}
\usepackage[hidelinks]{hyperref}

\begin{document}

\title{Covered: Browser Agents Cannot Tell a Clickable\\Element From an Occluded One}

\author{\IEEEauthorblockN{Bhargav Pedapudi}
\IEEEauthorblockA{Independent Researcher\\
San Francisco Bay Area, CA, USA\\
ORCID: 0009-0002-8523-8415\\
pedapudibhargav@gmail.com}}

\maketitle

\begin{abstract}
A browser agent is given a list of named interactive elements. When a consent
dialog covers the page, that list contains the covered elements exactly as it
contains the reachable ones. The agent selects the correct target, the click is
intercepted, and nothing in its observation explains the failure, so it repeats
the identical action until its step budget is exhausted. We audit the two tree
serialisations agents actually receive and find that neither expresses
occlusion on any of eight live sites carrying 25 to 141 covered elements, and
that the seven accessibility-tree properties Chrome defines for this purpose
are emitted on no node in any frame. We then evaluate the one deployed
detector: a paint-order rule recovers 0.84 to 0.95 of covered elements when
the blocker is in the main document and 0.00 to 0.17 when it is inside a
cross-origin iframe, missing 146 of 373 covered elements overall with zero
false positives, because paint order is computed per document while the
compositor is not. Hit testing closes both gaps in 3\,ms. Labelling occlusion
in the observation moves task success from 4/26 to 19/26 across eleven live
sites (McNemar exact two-sided $p = 0.00006$; seven sites better, none worse,
$p = 0.0156$) at fewer observation tokens, replicates on a second model, and
reproduces on a deterministic replay benchmark we release. A four-arm ablation
separates diagnosis from detection: a loop detector given a generic ``that
failed, try something else'' hint scores 4/17 against a flat baseline's 3/17,
so what recovers the agent is being told \emph{what} is wrong, not that
something is. We report five retracted claims of our own, including a
compression result that reversed sign under a proper paired test.
\end{abstract}

\begin{IEEEkeywords}
web agents, large language models, browser automation, accessibility tree,
occlusion, Chrome DevTools Protocol, reproducible benchmarks
\end{IEEEkeywords}

\section{Introduction}

An LLM browser agent acts through an observation of the page. In current
practice that observation is a flat list of interactive elements, each with a
role, an accessible name and a reference the agent can act on. The list is
produced either from the Chrome DevTools Protocol accessibility tree or from
Playwright's snapshot, and both are in wide deployment.

Neither can say that an element cannot be reached.

The consequence is a specific and reproducible failure. On a page where a
consent dialog intercepts pointer events, the agent is handed a list in which
the covered navigation links appear with correct names and correct references.
It selects the right one. The click times out. The observation on the next step
is unchanged, so the agent selects the same element again:

\begin{verbatim}
click('27') -> err:TimeoutError
click('27') -> err:TimeoutError
click('27') -> err:TimeoutError
... six identical attempts, task failed
\end{verbatim}

The agent did not choose badly. It was given a representation in which the
distinction it needed does not exist.

This paper makes that argument by measurement. Section~\ref{sec:audit} audits
what the deployed serialisations contain. Section~\ref{sec:paint} evaluates the
one shipped occlusion detector and shows that its method is blind to the most
common kind of full-page overlay. Section~\ref{sec:ablation} measures what
supplying the missing fact is worth, and separates it from the cheaper
alternative of detecting failure. Section~\ref{sec:cost} decomposes the cost of
a step. Section~\ref{sec:bench} describes the deterministic benchmark we
release. Section~\ref{sec:negative} reports what we got wrong.

\section{Related Work}

\textbf{Agents fail on consent banners, and it is documented without a cause.}
BrowserArena~\cite{browserarena} names pop-up banner removal as one of three
recurring open-web failure modes. Over 80 tasks on a single news site it
reports the banner detected in 100.00\% of Claude-3.7 runs and closed in only
7.5\%; Llama-4 detected 98.75\% and closed 17.72\%; o4-mini closed 17.81\% and
Gemini~2.5 closed 4.65\% (Appendix~K, Table~4). The failure is quantified; no
mechanism is given beyond the name of the task.

\textbf{Repeated ineffective actions are a studied failure class.}
MIRAGE-Bench~\cite{miragebench} constructs settings in which an agent is
persistently executing the same incorrect action, probing at the third and
sixth repetition (Appendix~A.6.1). A human audit of WebArena-Lite
trajectories~\cite{webarenalite} reviews 102 failed GPT-5.5 runs and reports
frequent scrolling loops. Neither attributes a loop to occlusion.

\textbf{The premise of this paper has been stated once, in passing.}
MIRAGE-Bench's pop-up setting appends pop-up elements to the end of the
accessibility tree, and its Figure~8 caption observes that the pop-up
``visually occludes part of the interface, but does not remove underlying
components from the accessibility tree''~\cite{miragebench}. That is the gap we
measure. It appears in a figure caption; it is not quantified, not fixed and
not ablated.

\textbf{Telling an agent to ignore a pop-up does not work.} Zhang et
al.~\cite{popupattack} achieve an 86\% attack success rate with adversarial
pop-ups, decreasing task success by 47\%, and report that basic defences such
as asking the agent to ignore pop-ups are ineffective. Their setting is
adversarial rather than benign consent, and they do not compare against an
observation-side remedy. Our Section~\ref{sec:ablation} supplies that
comparison for the benign case.

\textbf{Overlays measurably reduce agent success in industry testing.} A
practitioner study~\cite{merj} reports agent success rising from approximately
50\% to approximately 64\% on a full-screen consent banner between two
measurement rounds, and from approximately 32\% to approximately 41\% on a
transparent \texttt{opacity:0} overlay that swallows clicks. No mechanism is
offered.

\textbf{Occlusion is detectable, and one deployed system detects it.} Browser
Use marks a node \texttt{ignored\_by\_paint\_order} when its rectangle lies
inside the union of rectangles painted above it, using CDP paint
order~\cite{browseruse}. To our knowledge this has never been evaluated.
Section~\ref{sec:paint} evaluates it.

\textbf{Observation-format ablations vary modality, not reachability.}
Published comparisons set screenshots against accessibility trees, or
set-of-marks against text. Work on reducing the observation space includes
region-based reorganisation, which reports observation tokens falling from
6{,}437 to 3{,}671 and success improving by 2.3 percentage
points~\cite{region4web}. We are not aware of any work that varies whether the
agent is told which elements it can reach.

\section{What the Deployed Formats Contain}
\label{sec:audit}

\subsection{Method}

We load eight sites that present a blocking consent dialog and serialise each
page three ways: the CDP \texttt{Accessibility.getFullAXTree} call on the top
document; the same call issued per frame and concatenated, which is the
approach BrowserGym takes via \texttt{Page.getFrameTree}~\cite{browsergym}; and
Playwright's \texttt{aria\_snapshot(mode="ai")}, which is what Playwright MCP
serves to an agent. For each serialisation we ask two questions: does it
contain the control that dismisses the dialog, and does it state anywhere that
elements are covered.

For the second question we search each serialisation for any of the terms a
format would have to use to express occlusion, and we additionally enumerate
the accessibility-tree properties Chrome's protocol defines for this purpose:
\texttt{inertElement}, \texttt{inertSubtree}, \texttt{activeModalDialog},
\texttt{activeAriaModalDialog}, \texttt{ariaHiddenSubtree},
\texttt{notRendered} and \texttt{notVisible}.

\subsection{Result}

\begin{table}[t]
\caption{What each deployed serialisation contains, on eight live sites
carrying 25--141 covered elements each.}
\label{tab:audit}
\centering
\begin{tabular}{@{}lcc@{}}
\toprule
Serialisation & Dismiss control & Occlusion stated \\
\midrule
\texttt{getFullAXTree}, top document & 0 / 4 & \textbf{0 / 8} \\
\texttt{getFullAXTree}, merged per frame & 4 / 4 & \textbf{0 / 8} \\
\texttt{aria\_snapshot(mode="ai")} & 4 / 4 & \textbf{0 / 8} \\
\midrule
This work (hit tested) & 4 / 4 & 8 / 8 \\
\bottomrule
\end{tabular}
\end{table}

Table~\ref{tab:audit} gives the result. Two findings separate.

First, frame blindness belongs to one integration pattern, not to web agents
in general. A top-document-only \texttt{getFullAXTree} misses the dismiss
control on 4 of 4 sites where the consent panel renders inside a cross-origin
iframe. Issued per frame and merged, the same call finds it on 4 of 4. So does
Playwright's AI-mode snapshot, which emits the control with a frame-qualified
reference, for example \texttt{button "Accept and Continue" [ref=f4e22]}. A
naive integration is blind here; the tools are not.

Second, and this is the gap: \textbf{occlusion is expressed by none of the
three, on any of the eight sites.} Between 25 and 141 interactive elements per
page are covered and unclickable, and all three serialisations list them
exactly as they list reachable elements. Chrome's own vocabulary for this is
unused --- none of the seven properties above was emitted on any node, in any
frame, on any site.

One apparent exception was a false positive worth recording: a keyword search
matched the word ``interceptor'' in a news headline. We mention it because the
measurement is a search over text and such matches must be checked rather than
counted.

\section{Paint Order Misses the Overlays That Matter}
\label{sec:paint}

\subsection{Method}

Ground truth for ``would a click land here'' is
\texttt{document.elementFromPoint}, which reports what the compositor hits.
For each visible interactive element we hit-test its centre and treat the
element as covered when the returned node is neither the element nor an
ancestor nor a descendant of it.

Against that we evaluate the paint-order rule. We read real paint order from
\texttt{DOMSnapshot.captureSnapshot} with \texttt{includePaintOrder}, rather
than approximating it from document order, and mark an element covered when its
rectangle lies entirely inside the rectangle of a node painted above it, with
contributors below 0.8 opacity excluded as in the deployed
implementation~\cite{browseruse}.

An earlier version of this comparison approximated paint order from document
order and produced a different and misleading answer; we report the CDP-based
figures only.

\subsection{Result}

\begin{table}[t]
\caption{Paint-order recall against hit testing, eight live sites. The split
is by what is blocking, and it holds 8 of 8.}
\label{tab:paint}
\centering
\begin{tabular}{@{}lrrr@{}}
\toprule
Site & Covered & Found & Recall \\
\midrule
\multicolumn{4}{@{}l}{\emph{Blocker inside a cross-origin iframe}}\\
BBC             & 33 & 0  & \textbf{0.00} \\
Financial Times & 30 & 0  & \textbf{0.00} \\
Spiegel         & 27 & 1  & \textbf{0.04} \\
Sky Sports      & 42 & 7  & \textbf{0.17} \\
\midrule
\multicolumn{4}{@{}l}{\emph{Blocker in the main document}}\\
Le Monde        & 25 & 21 & 0.84 \\
Euronews        & 141 & 127 & 0.90 \\
CNN             & 36 & 34 & 0.94 \\
Irish Times     & 39 & 37 & 0.95 \\
\midrule
Total           & 373 & 227 & 0.61 \\
\bottomrule
\end{tabular}
\end{table}

Table~\ref{tab:paint} shows 146 of 373 covered elements missed, with
\emph{zero} false positives: the rule is conservative rather than noisy.

The failures are not spread evenly. Recall is near zero on four sites and
near one on four others, and the division is mechanistic. We classified each
page's blocking element and found, with no exceptions in eight sites, that the
four low-recall sites are exactly those whose consent dialog renders its panel
inside a cross-origin iframe.

The reason is that \textbf{paint order is computed per document}. When the
panel lives in a cross-origin iframe, the opaque rectangle that covers the page
is not a node in the top document's snapshot. Nothing in that document is
painted above the links, so nothing is marked covered --- while the compositor
places the frame on top and sends the click to it. A hit test does not have
this problem, because \texttt{elementFromPoint} reports what the compositor
actually hit, across document boundaries, and costs 3\,ms for a whole page.

This matters because the iframe case is the common case rather than the exotic
one. The major consent-management platforms render their panel in an iframe. In
an independent sweep of 42 further candidate sites, most of the sites we
rejected were rejected for exactly this reason --- blocked, with no dismiss
control reachable in the top document --- including pages with 59, 56, 181 and
252 covered elements.

\section{Does Labelling Occlusion Help?}
\label{sec:ablation}

\subsection{Setup}

We construct 26 tasks on 11 live sites. Each task's target is a navigation link
that sits behind a consent dialog, and success is a URL check, so nothing is
judged by a model. The task set was fixed and committed before the run.

All arms share the model, the action space and a six-step budget. Each task
gets a fresh browser context rather than cleared cookies, because
\texttt{clear\_cookies()} leaves \texttt{localStorage} intact and consent state
lives there; an earlier run without this inflated the effect when one arm
dismissed a dialog and its next tasks on the same site loaded with no dialog at
all.

The arms differ only in the observation:

\begin{itemize}
\item \textbf{A, flat.} The element list, as current practice produces it.
\item \textbf{B, structured.} The same list, with the overlay named, the
covered elements collapsed into a line of references, and the overlay's own
controls listed.
\item \textbf{B2, structured with four failure-driven additions}
(Section~\ref{sec:fixes}).
\item \textbf{C, flat with a specific hint} injected when an action repeats:
the hint names overlays and lists dismiss verbs.
\item \textbf{D, flat with a generic hint}: ``your last action was repeated and
had no effect. Something is preventing it from working. Try a different
approach.''
\end{itemize}

\subsection{Labelling works, on two models}

\begin{table}[t]
\caption{Occlusion labelling, 26 tasks on 11 live sites.}
\label{tab:main}
\centering
\begin{tabular}{@{}lrr@{}}
\toprule
 & \texttt{gpt-5-mini} & \texttt{gpt-5.4-mini} \\
\midrule
A, flat                   & 4 / 26  & 4 / 26 \\
B, structured             & 16 / 26 & \textbf{16 / 26} \\
B2, with fixes            & \textbf{19 / 26} & --- \\
\midrule
paired vs.\ flat          & $+15/-0$ & $+14/-2$ \\
McNemar exact two-sided   & $\mathbf{p = 0.00006}$ & $\mathbf{p = 0.00418}$ \\
by site                   & 7 better, 0 worse & 6 better, 1 worse \\
sign test                 & $p = 0.0156$ & $p = 0.125$ \\
\midrule
steps per task            & 5.23 $\to$ \textbf{3.38} & 5.23 $\to$ 3.77 \\
observation tokens        & 3{,}008 $\to$ \textbf{2{,}887} & --- \\
\bottomrule
\end{tabular}
\end{table}

Table~\ref{tab:main} gives the result. Eleven of the flipped tasks are the same
tasks under both models, which argues that the effect belongs to the page
rather than to one model's idiosyncrasies.

We report the site-level sign test alongside the task-level test because 26
tasks on 11 sites are not 26 independent trials: tasks on one site share an
overlay. Counting sites as the unit is the conservative reading, and it remains
significant for the first model.

The structured observation is \emph{smaller}. Collapsing 33 covered elements
into a line of references costs fewer tokens than listing all 33 as though they
worked, so the accuracy gain here does not trade against cost.

The two control sites, where the dialog is present but does not obstruct the
navigation, tie in both arms --- which is what a control should do, and which
rules out a general benefit from the richer format.

\subsection{Diagnosis, not detection}
\label{sec:diagnosis}

A loop detector is the obvious cheaper fix: the flat arm emits the same
timing-out click six times, so telling it ``that repeated and failed'' should
suffice. It does not.

\begin{table}[t]
\caption{Four arms, 17 tasks. Detecting the failure is not what helps.}
\label{tab:arms}
\centering
\begin{tabular}{@{}lrl@{}}
\toprule
Arm & Solved & vs.\ flat \\
\midrule
A, flat                        & 3 / 17  & --- \\
D, flat + \emph{generic} hint  & 4 / 17  & $+1/-0$, $p = 1.00$ \\
B, occlusion labelled          & 12 / 17 & $+10/-1$, $p = 0.0117$ \\
C, flat + \emph{specific} hint & 13 / 17 & $+10/-0$, $p = 0.0020$ \\
\bottomrule
\end{tabular}
\end{table}

Table~\ref{tab:arms} separates the two explanations. Given a hint that names
the failure mode and lists the remedy, the detector matches the structured
observation. Given a hint that says only that something is wrong, it is worth
one task out of seventeen and is statistically indistinguishable from silence.

The traces show the detector working and not helping:

\begin{verbatim}
click('27')  -> err:TimeoutError   correct target
click('27')  -> err:TimeoutError   repeat detected, hint fires
click('525') -> err:TimeoutError   varies, as instructed
click('27')  -> err:TimeoutError
click('21')  -> err:TimeoutError
click('593') -> err:TimeoutError   budget gone
\end{verbatim}

The agent obeys. It stops repeating and starts varying, and every variation
times out too, because everything it can see is behind the same overlay and
nothing in its observation mentions one. Varying blindly over a covered page is
not a recovery strategy.

Arm C's hint is an oracle: hand-written for the failure it was evaluated on,
correct for exactly one failure mode, and selectable only by a mechanism that
would already have to know what is wrong. It also pays a fixed tax. On the
tasks both C and B solve, B takes 2.33 steps and C takes 3.50, because a
detector cannot fire until the agent has already spent the actions that prove
the failure.

So the effect is carried by the diagnosis, not by the detection. Not ``you are
looping'' but ``you are covered''.

\subsection{Four fixes from four failures}
\label{sec:fixes}

We examined every task that failed in the structured arm and derived one change
from each cause: an imperative first line naming the control to activate; a
ranking over the overlay's controls; a parameter-free obstruction test, which
suppresses the announcement when the overlay covers nothing in a
\texttt{navigation} or \texttt{main} landmark; and a note when the dismiss
control is in an iframe and therefore absent from the listing.

Together these move 16/26 to 19/26 and cut steps by 19\%. Against flat the
improved arm is $+15/-0$, $p = 0.00006$, and by site 7 better and 0 worse,
$p = 0.0156$. \textbf{Against plain structure it is $+4/-1$, $p = 0.375$ ---
not significant.} The failure analysis improved the system; at $n = 26$ the
improvement itself is inside noise, and we state it as such.

The control ranking is correct on 2 of 6 sites using only generic structural
signals (a button rather than a link, \texttt{autofocus}, a declared default
action, area) and on 5 of 6 once an accept-word and trap-word list is added. The
vocabulary is doing the work, so we present the ranking as an engineering
default with a measured contribution rather than as a structural result.

\subsection{Wording decides whether the agent keeps trying}

One of the four fixes caused a regression, and it is the most transferable
result in this section.

On a site whose blocker is an advertisement iframe rather than a consent
dialog, the fourth fix emitted: ``No control in this document closes it.'' The
agent read a statement of absence as a statement of impossibility. It stopped
exploring and clicked one element five times, where the same arm without the
line had succeeded in four steps.

Reworded to state the same fact without closing the door --- ``its controls are
not listed here; the page behind it may still respond'' --- one of the two
affected tasks recovers.

\textbf{The channel that rescues an agent by naming what is covered can disable
it by implying nothing can be done.} An observation format is not a neutral
description of a page; its phrasing changes whether the agent continues to
search.

\section{The Cost of Not Knowing}
\label{sec:cost}

\begin{table}[t]
\caption{Where a step's wall clock goes, instrumented over 193 steps.}
\label{tab:cost}
\centering
\begin{tabular}{@{}lrrr@{}}
\toprule
Phase & Mean & Median & Share \\
\midrule
Observation extraction & 44\,ms  & 15\,ms  & 2.4\% \\
Model call             & 887\,ms & 833\,ms & 47.7\% \\
Browser action         & 928\,ms & 104\,ms & 49.9\% \\
\midrule
Step total             & 1{,}859\,ms & & \\
\bottomrule
\end{tabular}
\end{table}

Table~\ref{tab:cost} decomposes a step. Extraction is 2.4\% of it, so the
3\,ms hit test is not a cost worth optimising. The model call and the browser
action are roughly half each, which is consistent with a profiling study
reporting median delays of 4.7\,s for LLM inference and 6.6\,s for browser
action time~\cite{skim}.

The mean browser action time hides the structure that matters:

\begin{center}
\begin{tabular}{@{}lrr@{}}
\toprule
Action outcome & Mean & Median \\
\midrule
Lands ($n = 148$)       & 215\,ms   & \textbf{87\,ms} \\
Intercepted ($n = 45$)  & 3{,}390\,ms & \textbf{4{,}004\,ms} \\
\bottomrule
\end{tabular}
\end{center}

\textbf{An intercepted click costs four seconds, 46 times one that lands.} A
flat arm that spends six steps on a covered page therefore burns roughly 24
seconds in dead browser timeouts alone, on top of six model calls. Occlusion
blindness costs wall clock twice: in extra round trips, and in a four-second
stall on each blocked click.

Observation size is not the lever here. Over the range agents work in we
measured a 72$\times$ prompt range costing 26\% in latency; an earlier figure
of 11\% over a 32$\times$ range was an artifact of prompt caching, since the
sweep had sent five identical prompts per size.

\section{A Deterministic Benchmark}
\label{sec:bench}

Every number above was measured on live sites, and that is the weakest property
of the study. Five of eleven sites showed no overlay when re-probed minutes
after reporting 17 to 182 covered elements. Consent dialogs vary with
geography, experiment assignment and prior state, and nobody else can rerun a
measurement taken this way.

We therefore record each site's network traffic to HAR with the overlay present
and serve a fresh browser context entirely from that recording, with a request
the recording does not contain failing rather than reaching the network. A
recording is used only if it reproduces the condition under study: we require
the overlay to be present on replay and the covered-element count to be within
35\% of the live figure, and the experiment runs only the sites that pass.

Ten of eleven recordings reproduce their overlay, nine with an identical
covered-element count, including all four sites that had been unstable live.
The instability was in the network, not in the page.

The headline reproduces deterministically: 2/24 to 11/24, $+9/-0$, McNemar
exact two-sided $p = 0.00391$, at 2{,}925 to 2{,}525 observation tokens.
Magnitude is lower than live and the site-level sign test falls to 4 better and
0 worse, $p = 0.125$. Replay is also slower --- a step averages 3{,}141\,ms
against 1{,}859\,ms, with browser action time 2{,}300\,ms against 928\,ms ---
so the decomposition in Section~\ref{sec:cost} belongs to the live runs and the
two must not be pooled.

Building the fixture produced one instructive defect. The first recording
captured only each site's landing page, and the replayed experiment returned
0/17 in both arms, which reads as a clean refutation of the entire result. It
was the fixture: the success condition is a URL change, the destination pages
were never recorded, and an unrecorded navigation aborts, so an agent could
dismiss the overlay, click the correct link and still fail. Three things in the
data indicated this before any conclusion was drawn --- the improved arm's
observation was 2{,}011 tokens against the flat arm's 3{,}052, so the overlay
was being handled correctly; every task used the full budget in both arms; and
the recording contained exactly one document for the site under test. The
recorder now visits every target path during capture.

We release the code, the recordings and the raw results for every number in
this paper.

\section{Negative Results and Retractions}
\label{sec:negative}

We report these because they are what makes the positive result credible, and
because each was found by testing our own claim.

\textbf{Filtering the observation is harmful, not neutral.} We previously
reported compression as an empty direction on the basis of 14 tasks with no
paired test. At 40 paired tasks on a component benchmark, filtering the
observation to reachable, named, actionable elements cuts it by 80\% (median
1{,}607 to 197.5 characters) and costs 36/40 to 30/40, with 0 gained and 6
lost, McNemar $p = 0.03125$. Five of the six lost tasks go from one or two
steps to exhausting the budget, and they are a drawer, a collapsible
disclosure, a notification centre, a custom select and a text input --- controls
whose contents are hidden until revealed, which is exactly the set a
reachability filter removes.

\textbf{An attribution of that loss, which failed.} The appealing reading is a
symmetry with Section~\ref{sec:ablation}: labelling occlusion helps, deleting
it hurts, and the deployed practice of dropping covered nodes is thereby
indicted. One ablation refuted it. Keeping covered elements in the filtered
observation, with the other filters unchanged, recovered 0 of the 6 lost tasks
and scored slightly worse again (28/40). The harm belongs to the filtered
serialisation as a whole; which component carries it is unlocalised.

\textbf{Cross-frame traversal is not a contribution.} We implemented frame
traversal with frame-qualified references and reported it as closing a gap no
page observation could reach. Playwright's AI-mode snapshot already descends
into cross-origin iframes and already uses that reference scheme, and
BrowserGym already merges per-frame trees. The claim is withdrawn; what survives
is the narrower statement in Section~\ref{sec:audit} about a top-document-only
integration.

\textbf{Occlusion detection is not absent from deployed systems.} We claimed it
was. Browser Use ships it. Section~\ref{sec:paint} is the result of checking.

\textbf{A prevalence citation we cut.} A figure we had attributed to a consent
measurement study is not that study's statistic, and we removed the claim
rather than soften it. We are not aware of a prevalence figure for blocking
consent dialogs that we can stand behind, and we make no prevalence claim.

\section{Limitations}
\label{sec:limits}

The experiments cover one task family --- navigation, with success by URL check
--- and one overlay type, the consent dialog. Autocomplete menus, sticky
headers, interstitials and transparent click-swallowing overlays are not
measured. Twenty-six tasks on eleven sites is a small study, and the
improvement of the failure-driven additions over plain structure is not
significant at that size. Two models, from one family, are not model
independence. The paint-order comparison reimplements the deployed rule from its
source rather than running that system end to end, which could be unfair to it
in ways a reimplementation hides. One site reproduces its overlay on replay yet
collapses from 4/4 live to 0/4 replayed in both arms, and we have not diagnosed
it.

\section{Conclusion}

The observation formats browser agents are given cannot express the difference
between an element that can be clicked and one that cannot. The fact is
available from the DOM for 3\,ms and at negative token cost, the accessibility
protocol already defines a vocabulary for it that goes unused, and the one
deployed detector uses a per-document method blind to the overlays that occur
most. Supplying the fact moves success from 4/26 to 19/26 across eleven sites
and reproduces deterministically. What recovers the agent is being told what is
wrong rather than that something is: a loop detector with a generic hint is
worth one task in seventeen. And the same channel can disable the agent if it
states an absence as an impossibility.

\section*{Artifact Availability}

Code, the HAR recordings, and the raw results behind every number are at
\url{https://github.com/pedapudibhargav/glance-occlusion}.

\begin{thebibliography}{00}

\bibitem{browserarena}
S. Anupam, D. Brown, S. Li, E. Wong, H. Hassani, and O. Bastani,
``BrowserArena: Evaluating LLM agents on real-world web navigation tasks,''
arXiv:2510.02418, 2025.

\bibitem{miragebench}
W. Zhang, Y. Sun, P. Huang, J. Pu, H. Lin, and D. Song,
``MIRAGE-Bench: LLM agent is hallucinating and where to find them,''
arXiv:2507.21017, 2025.

\bibitem{webarenalite}
C. Gan, Z. He, Y. Tsujii, K. Kobayashi, H. Itoh, and K. Funakoshi,
``Auditing web agent evaluation on WebArena-Lite: Human review of outcomes and
trajectories,'' arXiv:2610.01491, 2026.

\bibitem{popupattack}
Y. Zhang, T. Yu, and D. Yang, ``Attacking vision-language computer agents via
pop-ups,'' in \emph{Proc. ACL}, 2025. arXiv:2411.02391.

\bibitem{merj}
Merj, ``How AI agents use websites, where they fail, and what to fix,''
2026. [Online]. Available: \url{https://merj.com/blog/agent-interaction-research}

\bibitem{browseruse}
Browser Use, ``\texttt{browser\_use/dom/serializer/paint\_order.py},''
source code. [Online]. Available: \url{https://github.com/browser-use/browser-use}

\bibitem{browsergym}
ServiceNow Research, ``\texttt{browsergym/core/observation.py},'' source code.
[Online]. Available: \url{https://github.com/ServiceNow/BrowserGym}

\bibitem{region4web}
D. Kwon and D. Lee, ``Rethinking observation space granularity for web
agents,'' arXiv:2605.07134, 2026.

\bibitem{skim}
M. Wong, K. Hsieh, S. Nath, and R. Netravali, ``Speculative execution for fast
and efficient web agents,'' arXiv:2605.16565, 2026.

\end{thebibliography}

\end{document}
